% Generated draft: requires mathematical and source review.
\documentclass[11pt]{article}
\usepackage[margin=25mm]{geometry}
\usepackage{fontspec}
\setmainfont{Latin Modern Roman}
\usepackage{amsmath,amssymb,mathtools,unicode-math}
\setmathfont{Latin Modern Math}
\usepackage{enumitem,needspace,xurl,hyperref}
\setlength{\parindent}{0pt}
\setlength{\parskip}{5pt}
\setlength{\emergencystretch}{3em}
\newcommand{\N}{\mathbb{N}}
\newcommand{\Z}{\mathbb{Z}}
\newcommand{\Q}{\mathbb{Q}}
\newcommand{\R}{\mathbb{R}}
\newcommand{\C}{\mathbb{C}}
\newcommand{\F}{\mathbb{F}}
\newcommand{\A}{\mathbb{A}}
\newcommand{\PP}{\mathbb{P}}
\newcommand{\E}{\mathbb{E}}
\newcommand{\eps}{\varepsilon}
\newcommand{\dd}{\,\mathrm{d}}
\begin{document}
\section*{UnsolvedMath ID: 30003686}
\textbf{Source problem code:} catalogue entry
\subsection{Definitions and mathematical statement}
Let \(\omega=\{0,1,2,\ldots\}\). For sets \(A,B\subseteq\omega\), enumeration reducibility \(A\leq_e B\) means that there is a computably enumerable collection \(W\) of pairs \((n,F)\), where \(n\in\omega\) and \(F\subseteq\omega\) is finite, such that
\[
 n\in A\quad\Longleftrightarrow\quad \exists F\subseteq B\;((n,F)\in W).
\]
An enumeration degree is an equivalence class under mutual enumeration reducibility; let \(\mathcal D_e\) denote the set of these degrees.

The supplied degree-theoretic question is whether
\[
 \forall \mathbf a\in\mathcal D_e,\qquad
 \bigl(\mathbf a\text{ is almost-total}\bigr)
 \Longrightarrow
 \bigl(\mathbf a\text{ is graph-cototal}\bigr).
\]
The record does not define almost-total or graph-cototal. These terms retain their intended meanings in the cited report; their precise definitions must be checked there before this formulation is fully self-contained.

The source presents the following topological question as equivalent. Let
\[
 H=[0,1]^\omega
\]
with the product topology. Using the standard convention, a space is \emph{punctiform} if it contains no nondegenerate continuum: every nonempty compact connected subspace has exactly one point. Does there exist a sequence \((P_n)_{n\in\omega}\) of subsets of \(H\), each given its subspace topology, such that
\[
 H=\bigcup_{n\in\omega}P_n
\]
and
\[
 \forall n\in\omega\;\forall K\subseteq P_n,\qquad
 \bigl(K\ne\varnothing\ \land\ K\text{ is compact and connected}\bigr)
 \Longrightarrow |K|=1?
\]
No closedness or measurability requirement on the pieces is supplied. The word ``decomposition'' does not specify whether the pieces must be disjoint. Under the stated punctiform convention, a countable cover can be made disjoint by replacing \(P_n\) with \(P_n\setminus\bigcup_{j<n}P_j\), since punctiformity is inherited by subspaces. Thus either reading gives the same existence question, allowing empty pieces.

The equivalence of the two questions is supplied by the record and is not independently established here. The intended punctiform convention and the specialized degree definitions require source review.
\subsection{Short English statement}
The question asks whether every almost-total enumeration degree is graph-cototal. The source gives an equivalent question: can the Hilbert cube be covered by countably many subspaces, each containing no compact connected subset with more than one point? The specialized degree definitions and the source's topological conventions require review.
\subsection{Sources}
\begin{itemize}
\item Local UnsolvedMath registry, supplied complete source record: permanent ID \(30003686\), source problem code \texttt{catalogue entry}, title ``Graph-Cototal Enumeration Degrees and Punctiform Decompositions.'' Collection: Oberwolfach Reports Open Problems. No individual proposer is identified.
\item Supplied source citation: \emph{Computability Theory (2019)}, Oberwolfach Reports, DOI \texttt{10.4171/owr/2018/1}, pp.\ 5--41. Source URL and supplied literature reference ``Oberwolfach report'': \url{https://doi.org/10.4171/owr/2018/1}. The background describes Oberwolfach Reports as documenting mathematical workshops and their open-problem discussions.
\item Literature checked at: \texttt{2026-08-21}. This is supplied dated curation, not independently verified. The supplied status is ``open.'' The supplied assessment says that the report explicitly presents the almost-total versus graph-cototal question and its punctiform-decomposition equivalent as open, and that no later resolution was located. The record calls this a dated literature triage derived from supplied OWR curation and requests independent verification of time-sensitive status and cited claims.
\item The supplied proposed year and year tag are \(2018\), while the citation labels the report ``Computability Theory (2019).'' This bibliographic distinction remains unresolved.
\end{itemize}
\end{document}
